\documentclass[11pt,a4paper]{article}
\usepackage[utf8]{inputenc}
\usepackage[T1]{fontenc}
\usepackage[margin=2cm, top=2.4cm]{geometry}
\usepackage{xcolor}
\usepackage{enumitem}
\usepackage{titlesec}
\usepackage{fancyhdr}
\usepackage{hyperref}

% Course color palette (copied from beamer preamble)
\definecolor{mlblue}{HTML}{0066CC}
\definecolor{mlpurple}{HTML}{3333B2}
\definecolor{mllavender}{HTML}{ADADE0}
\definecolor{mllavender2}{HTML}{C1C1E8}
\definecolor{mllavender3}{HTML}{CCCCEB}
\definecolor{mllavender4}{HTML}{D6D6EF}
\definecolor{mlorange}{HTML}{FF7F0E}
\definecolor{mlgreen}{HTML}{2CA02C}
\definecolor{mlred}{HTML}{D62728}
\definecolor{mlgray}{HTML}{7F7F7F}
\definecolor{lightgray}{HTML}{F0F0F0}
\definecolor{midgray}{HTML}{B4B4B4}

% Section heading colors
\titleformat{\section}{\large\bfseries\color{mlpurple}}{}{0em}{}[\vspace{-4pt}\textcolor{mllavender}{\rule{\linewidth}{1pt}}\vspace{2pt}]
\titleformat{\subsection}{\normalsize\bfseries\color{mlblue}}{}{0em}{\textbullet~}

\titlespacing*{\section}{0pt}{8pt}{4pt}
\titlespacing*{\subsection}{0pt}{5pt}{2pt}

% Header/footer
\pagestyle{fancy}
\fancyhf{}
\setlength{\headheight}{22pt}
\addtolength{\topmargin}{-10pt}
\renewcommand{\headrulewidth}{0.4pt}
\renewcommand{\headrule}{\hbox to\headwidth{\color{mllavender}\leaders\hrule height \headrulewidth\hfill}}
\fancyhead[L]{\small\color{mlpurple}\textbf{Cryptography \& Security} \textcolor{mlgray}{|} Lesson 02 \textcolor{mlgray}{|} Pre-Class Discovery Handout}
\fancyhead[R]{\small\color{mlgray}Prof.\ Dr.\ Joerg Osterrieder}
\fancyfoot[C]{\small\color{mlgray}\thepage}

% Compact list spacing
\setlist[itemize]{nosep, topsep=2pt, partopsep=0pt, leftmargin=1.4em}
\setlist[enumerate]{nosep, topsep=2pt, partopsep=0pt, leftmargin=1.6em}

% Activity box helper
\newcommand{\activitybox}[3]{%
  \noindent\colorbox{mllavender4}{\parbox{\dimexpr\linewidth-2\fboxsep}{%
    \textbf{\color{mlpurple}#1}\hfill\textcolor{mlgray}{\small #2}\\[2pt]#3%
  }}\par\vspace{4pt}%
}

% Thin ruled cell for fill-in tables
\newcommand{\fillcell}{\rule{2.2cm}{0.3pt}}

\begin{document}

% ──────────────────────────────────────────────────────────────────────────────
%  PAGE 1
% ──────────────────────────────────────────────────────────────────────────────

\section*{\color{mlpurple}How Does Cryptography Protect Your Money?}

\vspace{-2pt}
\noindent\textcolor{mlgray}{\small Work through the activities below \emph{before} our second session.
Bring your notes and answers to class.}
\vspace{4pt}

% ── Activity 1 ────────────────────────────────────────────────────────────────
\activitybox{Activity 1 --- The Caesar Cipher}{$\sim$7 min | pairs}{%
\textbf{Rule:} Shift each letter 3 positions forward (\texttt{A$\to$D}, \texttt{Z$\to$C}).

\vspace{2pt}
\textbf{Encrypt} \texttt{BLOCKCHAIN} with shift~3. Write your ciphertext here:
\quad\rule{5cm}{0.3pt}
\quad\textit{(answer: \texttt{EORFNFKDLQ})}

\vspace{3pt}
\textbf{Break it:} Your partner encrypts a 6-letter word with an unknown shift and hands you only
the ciphertext. Try all 25 shifts until you find the plaintext. Record the shift you found:
\quad\rule{2cm}{0.3pt}

\vspace{3pt}
\textbf{Kerckhoffs's principle:} The algorithm is public; only the key is secret.
\textit{Why does Bitcoin's open-source code not compromise its security?}
\vspace{4pt}
\noindent\rule{\linewidth}{0.3pt}
}

\vspace{2pt}

% ── Activity 2 ────────────────────────────────────────────────────────────────
\activitybox{Activity 2 --- Hash It Yourself}{$\sim$7 min | individual}{%
\textbf{ASCII values:} \texttt{a=97 \; b=98 \; c=99 \; s=115 \; t=116 \; C=67}
\quad\textbf{Hash rule:} sum all values, take \textbf{mod 256}.

\vspace{3pt}
\begin{center}
\renewcommand{\arraystretch}{1.35}
\begin{tabular}{|l|c|c|}
\hline
\textbf{Input} & \textbf{Sum of ASCII values} & \textbf{Hash (mod 256)} \\
\hline
\texttt{cat}  & \fillcell & \fillcell \\
\hline
\texttt{Cat}  & \fillcell & \fillcell \\
\hline
\texttt{cats} & \fillcell & \fillcell \\
\hline
\end{tabular}
\end{center}

\vspace{2pt}
\textit{Answers: \texttt{cat}$\to$56 \quad \texttt{Cat}$\to$26 \quad \texttt{cats}$\to$171}

\vspace{2pt}
Changing one character yields a completely different hash --- the \textbf{avalanche effect}.
Can you reconstruct the word from the number~56? Why not?
}

\vspace{2pt}

% ── Activity 3 ────────────────────────────────────────────────────────────────
\activitybox{Activity 3 --- The Envelope Analogy}{$\sim$7 min | groups of 3}{%
\textbf{Paper exercise:}
\begin{enumerate}
  \item Write a short secret, fold the paper, seal it in an envelope. Write your \textbf{name} on
        the outside (your \textcolor{mlblue}{\textbf{public address}}).
  \item Classmates seal messages to you in separate envelopes addressed by name.
  \item \textbf{Only you} open your envelopes.
\end{enumerate}

\textbf{Mapping:} public key $\approx$ name on envelope (anyone can send); private key $\approx$
ability to open (only you decrypt). \textit{Why is sharing your Bitcoin private key catastrophic?}
\vspace{4pt}
\noindent\rule{\linewidth}{0.3pt}
}

\newpage

% ──────────────────────────────────────────────────────────────────────────────
%  PAGE 2
% ──────────────────────────────────────────────────────────────────────────────

\section*{\color{mlpurple}Why Can't You Fake a Signature?}

\vspace{-2pt}
\noindent\textcolor{mlgray}{\small Digital signatures and key management are the last line of defence
for your coins.}
\vspace{4pt}

% ── Activity 4 ────────────────────────────────────────────────────────────────
\activitybox{Activity 4 --- Sign and Verify}{$\sim$8 min | pairs}{%
\begin{enumerate}
  \item Write \textit{``Send 2 BTC to Alice''} on a card and add your \textbf{unique personal mark}
        (distinctive doodle or symbol) --- your ``private key'' signature.
  \item Your neighbour examines it: genuine or fake?
  \item Your neighbour now tries to \textbf{forge} your mark on a blank card. How convincing?
\end{enumerate}

\textbf{ECDSA connection:} In Bitcoin every transaction is signed with the sender's
Elliptic Curve Digital Signature Algorithm (ECDSA) private key. The network verifies against the
public key --- mathematically impossible to forge. Changing even one satoshi invalidates the
signature.

\vspace{3pt}
\textit{Why must the message content be part of what is signed, not just the sender's identity?}
\vspace{4pt}
\noindent\rule{\linewidth}{0.3pt}
}

\vspace{2pt}

% ── Activity 5 ────────────────────────────────────────────────────────────────
\activitybox{Activity 5 --- The Key Management Challenge}{$\sim$8 min | individual $\to$ group}{%
\textbf{Scenario:} You control 5 Bitcoin wallets, each with a separate private key.

\vspace{3pt}
\begin{center}
\renewcommand{\arraystretch}{1.3}
\begin{tabular}{|p{4.5cm}|p{4.5cm}|p{3.6cm}|}
\hline
\textbf{Question} & \textbf{Your answer} & \textbf{Group answer} \\
\hline
How many keys to back up? & \rule{4.2cm}{0.3pt} & \rule{3.3cm}{0.3pt} \\[6pt]
\hline
What if one key is lost? & \rule{4.2cm}{0.3pt} & \rule{3.3cm}{0.3pt} \\[6pt]
\hline
What if a backup is found? & \rule{4.2cm}{0.3pt} & \rule{3.3cm}{0.3pt} \\[6pt]
\hline
\end{tabular}
\end{center}

\vspace{2pt}
\textbf{Solution --- HD Wallets \& Seed Phrases:} A Hierarchical Deterministic (HD) wallet derives
\emph{all} keys from one \textbf{seed phrase} (12--24 words). Back up once $\Rightarrow$ recover
everything. Lose the seed $\Rightarrow$ lose all coins. \textit{Order of words matters.}
}

\vspace{4pt}

% ── Reflection ────────────────────────────────────────────────────────────────
\section{Reflection Questions --- Bring Your Answers to Class}

\begin{enumerate}[leftmargin=1.8em, label=\textbf{\color{mlpurple}\arabic*.}]
  \item SHA-256 is a \emph{one-way} function: easy to compute, practically impossible to reverse.
        Why is this critical for Bitcoin mining and password storage?
        What would happen if SHA-256 were reversible?

  \vspace{7pt}\noindent\rule{\linewidth}{0.3pt}\vspace{2pt}\noindent\rule{\linewidth}{0.3pt}
  \vspace{3pt}

  \item Quantum computers could break the elliptic-curve cryptography protecting Bitcoin private
        keys (Shor's algorithm). What happens to Bitcoin then?
        What post-quantum solutions are being explored, and what trade-offs do they introduce?

  \vspace{7pt}\noindent\rule{\linewidth}{0.3pt}\vspace{2pt}\noindent\rule{\linewidth}{0.3pt}
  \vspace{3pt}

  \item \textit{``Not your keys, not your coins.''} Compare the trade-offs of
        \textbf{self-custody} (hardware wallet, seed phrase) versus funds on a
        \textbf{centralised exchange}. When is each option the wiser choice?

  \vspace{7pt}\noindent\rule{\linewidth}{0.3pt}\vspace{2pt}\noindent\rule{\linewidth}{0.3pt}
\end{enumerate}

\vspace{4pt}
\noindent\textcolor{mlgray}{\small\textit{Prepared by Prof.\ Dr.\ Joerg Osterrieder \,\textbullet\, Cryptography \& Security --- Lesson 02 \,\textbullet\, Pre-Class Discovery Handout}}

\end{document}
